\clearpage
\fancyhead[L]{\small BEMERKUNGEN VON SCHERING.}
\hypertarget{native-659-lower}{}
\begin{center}\textsc{Bemerkungen.}\end{center}

Dieser VI. Band der Werke von \textsc{Gauss} enthält die von \textsc{Gauss} der Öffentlichkeit übergebenen Abhandlungen, Aufsätze, Beobachtungen und Rezensionen, insofern sie den astronomischen Wissenschaften angehören.
Der Band ist dadurch schon umfangreicher geworden als jeder der übrigen Bände dieses Werkes.
Die im Nachlass befindlichen Handschriften zu diesem Gebiet, insbesondere eine Theorie der planetarischen Störungen und die Pallastafeln konnten also schwerlich noch aufgenommen werden, sondern müssen getrennt von diesem Band gedruckt werden.

Da die Briefe an \textsc{Zach} und \textsc{Lindenau} nicht mehr vorhanden sind, habe ich aus den von den Herausgebern der Monatlichen Correspondenz für Erd- und Himmelskunde veröffentlichten Zusammenstellungen der Nachrichten über die kleinen Planeten diejenigen Stellen hier abdrucken lassen, die \textsc{Gauss} betreffen und großenteils aus seinen Briefen entnommen waren, außerdem nur noch solche Stellen, die zum Verständnis der ersteren notwendig erschienen.
Die Ephemeriden habe ich nur in einzelnen wenigen Beispielen wiedergegeben.
Bei den Vergleichen von Elementen mit Beobachtungen habe ich Letztere beigefügt, wenn sie sich an anderen Stellen in der Monatlichen Correspondenz finden.
Für die dieser Zeitschrift entnommenen Beurteilungen von Schriften anderer Verfasser ist der Name des Berichterstatters durch Briefe von \textsc{Lindenau} an \textsc{Gauss} bestätigt, für die den Göttinger gelehrten Anzeigen entnommenen durch die Rechnungsbücher der Königlichen Gesellschaft der Wissenschaften.

Göttingen 1874 November. \hfill \textsc{Ernst Schering}.
